\documentclass[11pt]{article}

\usepackage[margin=1in]{geometry}
\usepackage{amsmath,amssymb,amsthm}
\usepackage{algorithm}
\usepackage{algpseudocode}
\usepackage{enumitem}
\usepackage{hyperref}
\usepackage{mathtools}

% ── theorem environments ──
\newtheorem{theorem}{Theorem}[section]
\newtheorem{lemma}[theorem]{Lemma}
\newtheorem{corollary}[theorem]{Corollary}
\newtheorem{proposition}[theorem]{Proposition}
\theoremstyle{definition}
\newtheorem{definition}[theorem]{Definition}
\newtheorem{remark}[theorem]{Remark}

% ── notation shortcuts ──
\newcommand{\St}{\mathcal{S}}          % state space
\newcommand{\Dim}{\mathcal{D}}         % dimension set
\newcommand{\Edg}[1]{E(#1)}            % edge set of dimension d
\newcommand{\sel}{\sigma}              % a state / selection
\newcommand{\clean}{\mathsf{C}}        % clean operator
\newcommand{\res}[2]{\rho(#1,#2)}      % resolvedVal
\newcommand{\ik}{\mathsf{irrKey}}      % irrelevance key
\newcommand{\ck}{\kappa}               % class key
\newcommand{\Reach}{\mathrm{Reach}}    % reachability predicate
\newcommand{\Vis}{\mathrm{Vis}}        % visible predicate
\newcommand{\Act}{\mathrm{Act}}        % activated predicate
\newcommand{\Dis}{\mathrm{Dis}}        % disabled predicate
\newcommand{\Irr}{\mathrm{Irr}}        % irrelevance predicate
\newcommand{\fp}{\mathrm{fp}}          % fingerprint
\newcommand{\Rep}{\mathrm{Rep}}        % class representative
\newcommand{\wildcard}{*}

\title{Decision Graph Formalization\\[4pt]
       \large Structure, Algorithms, and Reduction Theorems}
\author{Singularity Map Project}
\date{}

\begin{document}
\maketitle

\begin{abstract}
We formalize the decision graph underlying the Singularity Map as a
mathematical structure, define the algorithms computed over it
(reachability, dead-end detection, validation), and state four
\emph{reduction theorems} used by the graph walker.

The central question is: \textbf{do the reductions change the results?}
The baseline algorithm is a full depth-first search that branches on every
enabled edge at every dimension. The production implementation applies a
chain of four reductions---equivalence-class merging, irrelevance
wildcarding, superposition with deferred collapse, and variant
expansion---that dramatically shrink the search space. Each reduction
theorem asserts that the reduced traversal produces \emph{exactly the same
answers} for every algorithm (reachability, dead-end detection, edge
coverage, and all validation invariants) as the baseline DFS.

This document is intended as input for formal or semi-formal verification
of these reduction theorems.
\end{abstract}

\tableofcontents
\newpage

%==========================================================================
\section{Graph Structure}\label{sec:structure}
%==========================================================================

%--------------------------------------------------------------------------
\subsection{Decision Graph}
%--------------------------------------------------------------------------

\begin{definition}[Decision Graph]\label{def:graph}
A \emph{decision graph} is a tuple
$G = (\Dim,\, E,\, \mathcal{A},\, \mathcal{H},\, \mathcal{W},\,
      \mathcal{R},\, \mathcal{O})$
where:
\begin{enumerate}[label=(\roman*)]
  \item $\Dim = \{d_1,\dots,d_n\}$ is a finite ordered set of
        \textbf{dimensions}.
  \item For each $d \in \Dim$, $\Edg{d} = \{e_1,\dots,e_{m_d}\}$ is a
        nonempty finite set of \textbf{edges} (possible values).
  \item $\mathcal{A}$ assigns to each dimension $d$ an
        \textbf{activation predicate}
        $\Act_d : \St \to \{0,1\}$.
  \item $\mathcal{H}$ assigns to each dimension $d$ a
        \textbf{hide predicate} $H_d : \St \to \{0,1\}$.
  \item $\mathcal{W}$ assigns to each dimension--edge pair $(d,e)$ a
        \textbf{disable predicate}
        $\Dis_{d,e} : \St \to \{0,1\}$.
  \item $\mathcal{R}$ assigns to each dimension $d$ a
        \textbf{derivation rule}
        $R_d : \St \to \Edg{d} \cup \{\bot\}$.
  \item $\mathcal{O} = \{O_1,\dots,O_K\}$ is a set of
        \textbf{outcome predicates} $O_k : \St \to \{0,1\}$.
\end{enumerate}
\end{definition}

%--------------------------------------------------------------------------
\subsection{States}
%--------------------------------------------------------------------------

\begin{definition}[State]\label{def:state}
A \textbf{state} (or \emph{selection}) is a partial function
\[
  \sel : \Dim \rightharpoonup \textstyle\bigcup_{d\in\Dim} \Edg{d}
  \qquad\text{with}\quad
  \sel(d) \in \Edg{d} \cup \{\bot\}
  \;\;\forall\, d \in \Dim,
\]
where $\sel(d) = \bot$ means dimension $d$ is \emph{unanswered}.
The set of all states is $\St$.
\end{definition}

%--------------------------------------------------------------------------
\subsection{Visibility and Activation}
%--------------------------------------------------------------------------

\begin{definition}[Visibility]\label{def:vis}
A dimension $d$ is \textbf{visible} in state $\sel$ iff it is activated
and not hidden:
\[
  \Vis(d,\sel) \;\iff\; \Act_d(\sel) = 1 \;\wedge\; H_d(\sel) = 0.
\]
Activation and hiding are evaluated using the \emph{resolved} values of
$\sel$ (Definition~\ref{def:resolved}).
\end{definition}

%--------------------------------------------------------------------------
\subsection{Resolved Values (Derivation)}
%--------------------------------------------------------------------------

\begin{definition}[Resolved Value]\label{def:resolved}
The \textbf{resolved value} of dimension $d$ in state $\sel$ is
\[
  \res{\sel}{d} \;=\;
  \begin{cases}
    R_d(\sel) & \text{if } R_d(\sel) \neq \bot,\\
    \sel(d)   & \text{otherwise}.
  \end{cases}
\]
The derivation rule $R_d$ may reference $\res{\sel}{d'}$ for other
dimensions $d'$; cycles are broken by returning the raw value $\sel(d)$
when a dimension is encountered while already being computed.
\end{definition}

%--------------------------------------------------------------------------
\subsection{Condition Matching}\label{sec:conditions}
%--------------------------------------------------------------------------

Activation, hiding, and disabling predicates are defined by
\emph{conditions}---conjunctions of per-dimension constraints. A single
condition $\varphi$ is a partial function mapping dimensions to
\emph{constraint types}:

\begin{definition}[Condition Types]\label{def:condtypes}
Each constraint on dimension $d$ in condition $\varphi$ takes one of five
forms:
\begin{center}
\begin{tabular}{lll}
\textbf{Tag} & \textbf{Notation} & \textbf{Matches $\sel$ iff} \\
\hline
Presence  & $\varphi(d) = \mathsf{true}$   & $\res{\sel}{d} \neq \bot$ \\
Absence   & $\varphi(d) = \mathsf{false}$  & $\res{\sel}{d} = \bot$ \\
Inclusion & $\varphi(d) = V \subseteq \Edg{d}$ & $\res{\sel}{d} \in V$ \\
Exclusion & $\varphi(d) = \neg V$            & $\res{\sel}{d} \neq \bot \;\Rightarrow\; \res{\sel}{d} \notin V$ \\
Required-exclusion & $\varphi(d) = \neg^! V$ & $\res{\sel}{d} \neq \bot \;\wedge\; \res{\sel}{d} \notin V$ \\
\end{tabular}
\end{center}
A condition $\varphi$ \textbf{matches} state $\sel$ iff every constraint
in $\varphi$ is satisfied.
\end{definition}

\begin{definition}[Predicates from Conditions]\label{def:pred}
\begin{itemize}
  \item $\Act_d(\sel) = 1$ iff $d$ has no activation conditions, or at
        least one activation condition matches $\sel$.
  \item $H_d(\sel) = 1$ iff some hide condition of $d$ matches $\sel$.
  \item $\Dis_{d,e}(\sel) = 1$ iff some \texttt{disabledWhen} condition
        of edge $e$ matches $\sel$, \emph{or} edge $e$ has a
        \texttt{requires} clause and no requirement condition matches.
\end{itemize}
\end{definition}

\begin{remark}[Direct Optimization]\label{rem:direct}
When every dimension referenced by a compiled condition is
\emph{non-derived} (has no derivation rule), the condition is flagged
\texttt{\_direct} and matched using raw state values $\sel(d)$
instead of $\res{\sel}{d}$, avoiding derivation computation. This is a
pure optimization; correctness follows because
$\res{\sel}{d} = \sel(d)$ when $R_d \equiv \bot$.
\end{remark}

%--------------------------------------------------------------------------
\subsection{Locking and the Clean Operator}\label{sec:clean}
%--------------------------------------------------------------------------

\begin{definition}[Locked Dimension]\label{def:locked}
A visible, non-derived dimension $d$ is \textbf{locked} in state $\sel$
if exactly one edge of $d$ is enabled:
$|\{e \in \Edg{d} : \Dis_{d,e}(\sel)=0\}| = 1$. The unique enabled edge
is $\ell_d(\sel)$.
\end{definition}

\begin{definition}[Clean Operator]\label{def:clean}
The \textbf{clean operator} $\clean : \St \to \St$ is the limit of the
iteration $\sel^{(0)} = \sel$,
\[
  \sel^{(i+1)}(d) =
  \begin{cases}
    \bot
      & \text{if } \sel^{(i)}(d) \neq \bot \text{ and }
        \Dis_{d,\sel^{(i)}(d)}(\sel^{(i)}) = 1,\\
    \sel^{(i)}(d) & \text{otherwise},
  \end{cases}
\]
run for at most $P$ passes (currently $P=5$) or until a fixpoint is
reached.

$\clean$ only removes disabled choices; it never writes new values.
Locked dimensions (where only one edge is enabled) are \emph{not}
auto-committed into $\sel$ — they are detected on demand via $\ell_d$
and surfaced to the user as a single-option confirmation step, which
the user commits through the normal push flow.
\end{definition}

\begin{lemma}[Termination]\label{lem:clean-term}
$\clean(\sel)$ terminates in at most $P \cdot |\Dim|$ steps, since each
pass inspects every dimension once and terminates early when no change
occurs.
\end{lemma}

%--------------------------------------------------------------------------
\subsection{Transitions and Plays}\label{sec:transitions}
%--------------------------------------------------------------------------

\begin{definition}[Transition]\label{def:trans}
From state $\sel$, choosing dimension $d$ with edge $e$ (where
$\Vis(d,\sel)=1$, $\sel(d)=\bot$, and $\Dis_{d,e}(\sel)=0$) yields
successor state:
\[
  \sel' = \clean(\sel[d \mapsto e])
\]
where $\sel[d \mapsto e]$ denotes $\sel$ with $\sel(d) \coloneqq e$.
\end{definition}

\begin{definition}[Play and Terminal]\label{def:play}
A \textbf{play} is a sequence of transitions
$\sel_0 \to \sel_1 \to \cdots \to \sel_T$ starting from the initial state
$\sel_0 = \clean(\lambda d.\,\bot)$.

State $\sel$ is \textbf{terminal} if no dimension is both visible and
unanswered:
\[
  \mathrm{Terminal}(\sel) \;\iff\;
  \nexists\, d \in \Dim :\;
  \Vis(d,\sel) = 1 \;\wedge\; \sel(d) = \bot.
\]

State $\sel$ is a \textbf{dead end} if some dimension is visible and
unanswered but every edge of that dimension is disabled.
\end{definition}

%--------------------------------------------------------------------------
\subsection{Outcome Matching}\label{sec:outcomes}
%--------------------------------------------------------------------------

\begin{definition}[Outcome Predicate]\label{def:outcome}
Each outcome $O_k$ is defined by a disjunction of \emph{reachable
conditions} over the \emph{resolved state}
$\hat\sel(d) = \res{\sel}{d}$:
\[
  O_k(\sel) = 1 \;\iff\;
  \exists\, \text{condition } c \in \mathcal{C}_k :\;
    \forall\,(d,V) \in c,\; \hat\sel(d) \in V
    \;\wedge\;
    \forall\,(d,V) \in c.\text{\_not},\; \hat\sel(d) \notin V.
\]
Some outcomes have a \textbf{primary dimension} $p_k$, producing
\emph{variants} $O_{k,v}$ that additionally require
$\hat\sel(p_k) = v$.
\end{definition}


%==========================================================================
\section{Algorithms}\label{sec:algorithms}
%==========================================================================

%--------------------------------------------------------------------------
\subsection{Baseline DFS}\label{sec:baseline}
%--------------------------------------------------------------------------

\begin{definition}[Next-Dimension Selection]\label{def:next}
Given state $\sel$, define $\mathsf{next}(\sel)$ as the first dimension
in the fixed ordering $d_1,\dots,d_n$ such that $\Vis(d,\sel)=1$ and
$\sel(d)=\bot$, respecting priority-based activation rules that may block
lower-priority dimensions when a higher-priority one is active but
unanswered. If no such dimension exists, $\mathsf{next}(\sel) = \bot$.
\end{definition}

\begin{algorithm}[H]
\caption{Baseline DFS for Reachability}\label{alg:baseline}
\begin{algorithmic}[1]
\Procedure{BaselineDFS}{$\sel$}
  \If{$\mathrm{Terminal}(\sel)$}
    \State \Return $M(\sel) \gets \sum_{k:\,O_k(\sel)=1} 2^k$
  \EndIf
  \State $d \gets \mathsf{next}(\sel)$
  \If{$d = \bot$} \Return $0$ \Comment{dead end}
  \EndIf
  \State $\mathit{mask} \gets 0$
  \ForAll{$e \in \Edg{d}$ with $\Dis_{d,e}(\sel)=0$}
    \State $\sel' \gets \clean(\sel[d \mapsto e])$
    \State $\mathit{mask} \gets \mathit{mask} \mathbin{|} \Call{BaselineDFS}{\sel'}$
  \EndFor
  \State \Return $\mathit{mask}$
\EndProcedure
\end{algorithmic}
\end{algorithm}

\begin{definition}[Reachability]\label{def:reach}
Outcome $k$ is \textbf{reachable} from state $\sel$ iff:
\[
  \Reach(\sel,k) \;\iff\;
  \text{bit } k \text{ is set in }
  \textsf{BaselineDFS}(\sel).
\]
Equivalently, there exists a play from $\sel$ reaching a terminal $\sel_T$
with $O_k(\sel_T) = 1$.
\end{definition}

%--------------------------------------------------------------------------
\subsection{Reachability Precomputation}\label{sec:precomp}
%--------------------------------------------------------------------------

In practice, the DFS is run once from $\sel_0$ with all $K$ outcome
matchers simultaneously. The result is a map
\[
  \mathsf{ReachMap} : \text{Key} \to \{0,\dots,2^K-1\}
\]
keyed by $\ik(\sel)$ (Definition~\ref{def:irrkey} below), storing the
OR-combined bitmask of reachable outcomes from any state mapping to that
key. Separate files are written per outcome (or variant): outcome $k$'s
file contains the set $\{\ik(\sel) : \text{bit } k \text{ set in }
\mathsf{ReachMap}(\ik(\sel))\}$.

%--------------------------------------------------------------------------
\subsection{Client-Side Lookup}\label{sec:client}
%--------------------------------------------------------------------------

At runtime, given the user's current choice stack and a candidate
move $(d,e)$:

\begin{enumerate}
  \item Take the current committed state $\sel$ from the stack.
  \item Apply $\clean(\sel[d \mapsto e])$ to get the hypothetical next
        state $\sel'$.
  \item Compute $\ik(\sel')$ and check membership in the preloaded
        $\mathsf{ReachSet}_k$.
\end{enumerate}

Because $\clean$ never auto-commits locked values, the committed client
state and the server-side DFS state inhabit the same key space, so
$\ik$ agreement is immediate.

%--------------------------------------------------------------------------
\subsection{Graph Walk}\label{sec:walk}
%--------------------------------------------------------------------------

The \textsf{walk} procedure traverses the full state space using the
superposition DFS (Algorithm~\ref{alg:superdfs}) and classifies every
leaf.

\begin{definition}[Walk Leaf Classification]\label{def:leafclass}
A leaf reached by the walk is classified as one of:
\begin{enumerate}[label=(\roman*)]
  \item \textbf{Terminal} --- $\mathrm{isTerminal}(\sel) = 1$
        (at least one outcome matched).
  \item \textbf{Dead end} --- $\mathsf{next}(\sel) = \bot$
        (unanswered visible dimensions may exist but no outcome matched),
        or $\mathsf{next}(\sel) = d$ but every edge of $d$ is disabled.
  \item \textbf{Boundary} --- $\mathsf{next}(\sel) = d$ and $d$ belongs
        to a specified \emph{exclude} set (the walk stops before entering
        those dimensions).
\end{enumerate}
\end{definition}

The walk also collects \textbf{edge coverage}: the set
$\{(d,e) : \exists\,\text{visited } \sel \text{ with } \sel(d)=e\}$,
including values produced by derivation.

%--------------------------------------------------------------------------
\subsection{Validation Algorithms}\label{sec:validate}
%--------------------------------------------------------------------------

The validator runs three phases over the graph. Phase~1 is purely static;
Phase~2 piggybacks on the \textsf{walk} DFS; Phase~3 concerns narrative
content and is omitted here (it is domain-specific and does not interact
with graph structure).

\subsubsection{Phase 1: Static Analysis}\label{sec:static}

These checks inspect the graph definition without traversal.

\begin{definition}[Reference Integrity]\label{def:refint}
For every condition $\varphi$ appearing in any \texttt{activateWhen},
\texttt{hideWhen}, \texttt{disabledWhen}, \texttt{requires}, or
\texttt{effects.when} rule: every dimension $d$ referenced
by $\varphi$ must exist ($d \in \Dim$), and every value referenced
must be a valid edge ($v \in \Edg{d}$). The same holds for
\texttt{effects.set} / \texttt{setFlavor} write maps,
module \texttt{exitPlan} \texttt{set} maps, and outcome template
conditions.
\end{definition}

\begin{definition}[Dead Edge Detection]\label{def:deaded}
An edge $e$ of dimension $d$ is \textbf{statically dead} if $e$ has
both a \texttt{requires} clause and a \texttt{disabledWhen} clause, and
for every requirement condition $r$ in \texttt{requires}, there exists a
disable condition $w$ in \texttt{disabledWhen} such that $w$ is
\emph{logically implied} by $r$:
\[
  \forall\, r \in \mathsf{requires}(e),\;
  \exists\, w \in \mathsf{disabledWhen}(e) :\;
  \forall\, (d_w, V_w) \in w,\;\;
    d_w \in \mathrm{dom}(r) \;\wedge\;
    r(d_w) \subseteq V_w.
\]
Such an edge can never be enabled: whenever its requirements are met, a
disable condition is also triggered.
\end{definition}

\subsubsection{Phase 2: DFS Invariant Checks}\label{sec:dfsinvariants}

At every state $\sel$ visited during the walk, the following invariants
are checked.

\begin{definition}[Dead End]\label{def:deadend}
A \textbf{dead end} is a non-terminal leaf: $\mathsf{next}(\sel) = \bot$
and no outcome predicate matches, \emph{or} $\mathsf{next}(\sel) = d$
with $\Edg{d}$ nonempty but
$\forall\, e \in \Edg{d}:\;\Dis_{d,e}(\sel)=1$.
A well-formed graph should have zero dead ends.
\end{definition}

\begin{definition}[Re-Derived Dead End]\label{def:rederived}
At a dead end $\sel$, a \textbf{re-derived dead end} occurs when some
dimension $d$ has a raw value $\sel(d)$ different from its resolved value
$\res{\sel}{d}$, and replacing raw values with resolved values reveals a
new visible unanswered dimension. This indicates a derivation rule that
hides a question that should be reachable.
\end{definition}

\begin{definition}[Stuck Node]\label{def:stuck}
A \textbf{stuck node} at state $\sel$ is a visible, non-derived,
non-locked, unanswered dimension $d$ for which every edge is disabled:
\[
  \Vis(d,\sel) = 1 \;\wedge\;
  \sel(d) = \bot \;\wedge\;
  \forall\, e \in \Edg{d}:\; \Dis_{d,e}(\sel) = 1.
\]
Unlike a dead end (which concerns only $\mathsf{next}(\sel)$), a stuck
node can exist at any dimension, even if a different dimension is next.
\end{definition}

\begin{definition}[Ambiguous Outcome]\label{def:ambig}
A state $\sel$ is \textbf{ambiguous} if it matches more than one outcome
template: $|\{k : O_k(\sel) = 1\}| > 1$. Well-formed outcomes should
be mutually exclusive at terminals.
\end{definition}

\begin{definition}[Single-Option Node]\label{def:singleopt}
A \textbf{single-option node} at state $\sel$ is a visible, unanswered,
\emph{non-locked} dimension $d$ where exactly one edge is enabled but
$d$ is not detected as locked by the engine. This is a UI concern: the
user sees a ``choice'' with only one option.
\end{definition}

\begin{definition}[Stack Integrity]\label{def:stackint}
At state $\sel$ with stack $\mathit{stk}$, the \textbf{stack integrity}
invariant holds iff the answered dimensions in \textsf{displayOrder}
appear in the same relative order as the non-derived entries in
$\mathit{stk}$. Violation indicates a dimension ordering inconsistency
between the engine's display logic and the traversal path.
\end{definition}

\subsubsection{Phase 2: Edge Coverage}\label{sec:coverage}

After the walk, \textbf{edge coverage} is the fraction of
$(d,e) \in \Dim \times E$ pairs that appear in at least one visited
state (including derived values). Unreached non-derived edges indicate
potential dead content or over-constrained activation rules.

\begin{definition}[Edge Coverage]\label{def:edgecov}
\[
  \mathrm{Coverage} =
  \frac{|\{(d,e) : \exists\,\text{visited } \sel,\;
         \sel(d) = e \text{ or } \res{\sel}{d} = e\}|}
       {|\{(d,e) : d \in \Dim,\; e \in \Edg{d}\}|}.
\]
\end{definition}


%==========================================================================
\section{Reductions}\label{sec:reductions}
%==========================================================================

The baseline DFS (Algorithm~\ref{alg:baseline}) is correct by
construction but explores up to $\prod_{d} |\Edg{d}|$ states---far too
many for practical use. The production traversal applies a chain of four
\textbf{reductions} that shrink the search space while (we claim)
preserving all results.

The reduction chain is:
\[
  \underbrace{\textsf{BaselineDFS}}_{\text{Section~\ref{sec:baseline}}}
  \;\xrightarrow{\text{R1}}\;
  \underbrace{\textsf{ClassDFS}}_{\text{\S\ref{sec:classes}}}
  \;\xrightarrow{\text{R2}}\;
  \underbrace{\textsf{ClassDFS} + \ik\text{-memo}}_{\text{\S\ref{sec:irrelevance}}}
  \;\xrightarrow{\text{R3}}\;
  \underbrace{\textsf{SuperDFS}}_{\text{\S\ref{sec:superposition}}}
  \;\xrightarrow{\text{R4}}\;
  \underbrace{\textsf{SuperDFS} + \text{Expand}}_{\text{\S\ref{sec:variants}}}
\]

\medskip
\noindent
\textbf{What ``preserves all results'' means.}
Let $\mathcal{A}$ denote the set of algorithms from
Section~\ref{sec:algorithms}: reachability (Def.~\ref{def:reach}),
dead-end detection (Def.~\ref{def:deadend}), edge coverage
(Def.~\ref{def:edgecov}), and the validation invariants of
Section~\ref{sec:validate} (stuck nodes, ambiguous outcomes, etc.).
For each reduction $\textsf{X} \xrightarrow{\text{R}} \textsf{Y}$ we
want to show:
\[
  \forall\, A \in \mathcal{A},\;
  \forall\, \sel \in \St :\;
  A_{\textsf{X}}(\sel) = A_{\textsf{Y}}(\sel).
\]
That is, the reduced traversal reports the same reachability bitmasks,
detects the same dead ends, visits the same edge-coverage set (up to
equivalence class), and triggers the same validation violations as the
baseline.

\medskip
\noindent
Each subsection below defines one reduction step and states the
corresponding theorem. Proofs are the target of formal verification.

%--------------------------------------------------------------------------
\subsection{Equivalence Classes}\label{sec:classes}
%--------------------------------------------------------------------------

\begin{definition}[Behavioral Equivalence]\label{def:equiv}
For dimension $d$, edges $e_1, e_2 \in \Edg{d}$ are
\textbf{behaviorally equivalent}, written $e_1 \sim_d e_2$, if for every
state $\sel$ with $\sel(d) \in \{e_1, e_2\}$, replacing $\sel(d)$ from
$e_1$ to $e_2$ (or vice versa) does not change:
\begin{enumerate}[label=(\alph*)]
  \item $\Vis(d',\sel)$ for any $d' \neq d$,
  \item $\Dis_{d',e'}(\sel)$ for any $(d',e')$ with $d' \neq d$,
  \item $\res{\sel}{d'}$ for any derived dimension $d'$.
\end{enumerate}
\end{definition}

This is an equivalence relation. It partitions $\Edg{d}$ into
\textbf{equivalence classes} $[e]_d$.

\begin{definition}[Class Key]\label{def:classkey}
The \textbf{class key} of state $\sel$ is:
\[
  \ck(\sel) = \bigl(c_1,\dots,c_n\bigr)
  \quad\text{where}\quad
  c_i =
  \begin{cases}
    [\res{\sel}{d_i}]_{d_i}
      & \text{if } d_i \text{ is derived},\\
    [\sel(d_i)]_{d_i}
      & \text{if } \sel(d_i) \neq \bot,\\
    \mathsf{U}
      & \text{if } \sel(d_i) = \bot.
  \end{cases}
\]
\end{definition}

\begin{definition}[Class Representative]\label{def:rep}
For each class $[e]_d$, fix a canonical \textbf{representative}
$\Rep_d([e]_d) \in [e]_d$.
\end{definition}

\begin{algorithm}[H]
\caption{Class-Merged DFS}\label{alg:classdfs}
\begin{algorithmic}[1]
\Procedure{ClassDFS}{$\sel$}
  \If{$\mathrm{Terminal}(\sel)$}
    \State \Return $M(\sel)$
  \EndIf
  \State $d \gets \mathsf{next}(\sel)$; \textbf{if} $d=\bot$ \textbf{return} $0$
  \State $\mathit{mask} \gets 0$
  \State $\mathit{seen} \gets \emptyset$
  \ForAll{$e \in \Edg{d}$ with $\Dis_{d,e}(\sel)=0$}
    \If{$[e]_d \notin \mathit{seen}$}
      \State $\mathit{seen} \gets \mathit{seen} \cup \{[e]_d\}$
      \State $\sel' \gets \clean(\sel[d \mapsto \Rep_d([e]_d)])$
      \State $\mathit{mask} \gets \mathit{mask} \mathbin{|} \Call{ClassDFS}{\sel'}$
    \EndIf
  \EndFor
  \State \Return $\mathit{mask}$
\EndProcedure
\end{algorithmic}
\end{algorithm}

\begin{theorem}[Reduction R1: Equivalence Class Correctness]\label{thm:class}
For all states $\sel$ and outcomes $k$:
\[
  \Reach_{\textsf{Baseline}}(\sel,k) = \Reach_{\textsf{ClassDFS}}(\sel,k).
\]
More broadly, the class-merged DFS detects the same dead ends (by class
key), reports the same validation violations, and covers the same edges
(up to within-class equivalence) as the baseline. Exploring only class
representatives at each branching point does not change any algorithm
result.
\end{theorem}

\begin{remark}[Computation]
Equivalence classes are computed by iterative signature-based clustering.
For each dimension $d$ and edge $e$, a \emph{signature} string encodes how
$e$ interacts with every condition, edge constraint, and derivation rule
in the graph. Edges with identical signatures are placed in the same
class. The process iterates (up to 20 rounds) until stable, then applies a
transitive collapse: if every ``reader'' of dimension $d$ has been
collapsed to a single class, dimension $d$ itself is collapsed.
\end{remark}

%--------------------------------------------------------------------------
\subsection{Irrelevance Wildcarding}\label{sec:irrelevance}
%--------------------------------------------------------------------------

\begin{definition}[Irrelevance]\label{def:irr}
Dimension $d$ is \textbf{irrelevant} at state $\sel$ (with class key
$\ck(\sel)$) if neither of the following holds:
\begin{enumerate}[label=(\alph*)]
  \item \textbf{Direct reader dependency:} There exists a non-derived
        dimension $d'$ that references $d$ in its activation, hiding, or
        edge conditions, $d'$ is unanswered ($\sel(d')=\bot$), and $d'$
        could become visible from $\sel$.
  \item \textbf{Derived reader dependency:} There exists a derived
        dimension $d'$ that reads $d$ via its derivation rule, changing
        $d$'s value within its class would change $\res{\sel}{d'}$, and
        $d'$ is itself \emph{not} irrelevant (recursive check, with
        cycles broken by treating revisited dimensions as irrelevant).
\end{enumerate}
We write $\Irr(d,\sel)$ when $d$ is irrelevant at $\sel$.
\end{definition}

\begin{definition}[Irrelevance Key]\label{def:irrkey}
The \textbf{irrelevance key} of state $\sel$ is:
\[
  \ik(\sel) = \bigl(k_1,\dots,k_n\bigr)
  \quad\text{where}\quad
  k_i =
  \begin{cases}
    \wildcard & \text{if } \Irr(d_i,\sel)
                \text{ and } d_i \text{ can be wildcarded},\\
    c_i       & \text{otherwise (class id from } \ck(\sel)\text{)}.
  \end{cases}
\]
A dimension can be wildcarded if: its value is defined, or it has no node,
or it is derived, or it is not currently visible.
\end{definition}

\begin{theorem}[Reduction R2: Irrelevance Correctness]\label{thm:irr}
For all states $\sel_1, \sel_2$ and outcomes $k$:
\[
  \ik(\sel_1) = \ik(\sel_2)
  \;\Longrightarrow\;
  \Reach(\sel_1,k) = \Reach(\sel_2,k).
\]
Consequently, memoizing by $\ik$ instead of by full state is safe: the DFS
can skip a state whose $\ik$ has already been visited, and all algorithms
(reachability, dead-end detection, validation) produce the same results.
\end{theorem}

\begin{remark}
Irrelevance enables aggressive memoization: the DFS stores results keyed
by $\ik(\sel)$ rather than the full state, collapsing the effective state
space. In the Singularity Map, the optimized walk visits
${\sim}216\text{k}$ distinct states; after variant and superposition
expansion at recording time (Sections~\ref{sec:variants},
\ref{sec:superexpand}), the final $\mathsf{ReachMap}$ contains
${\sim}243\text{k}$ distinct irrelevance keys covering the full 76-node,
13-outcome-family graph.
\end{remark}

%--------------------------------------------------------------------------
\subsection{Superposition}\label{sec:superposition}
%--------------------------------------------------------------------------

\begin{definition}[State Fingerprint]\label{def:fingerprint}
The \textbf{fingerprint} of state $\sel$ (at class key $\ck(\sel)$)
is:
\[
  \fp(\sel) = \bigl\langle\,
    (d', \{[e']_{d'} : e' \in \Edg{d'},\;\Dis_{d',e'}(\sel)=0\})
    \;\big|\;
    d' \in \Dim,\;\Vis(d',\sel)=1,\;\sel(d')=\bot,\;
    \Edg{d'} \neq \emptyset
  \,\bigr\rangle
\]
ordered by dimension ordering.
\end{definition}

\begin{definition}[Superimposability]\label{def:super}
At state $\sel$ with next dimension $d = \mathsf{next}(\sel)$, let
$C_1,\dots,C_m$ be the distinct equivalence classes among enabled edges
of $d$. Dimension $d$ is \textbf{superimposable} at $\sel$ if $m \geq 2$
and:
\[
  \forall\, i,j \in \{1,\dots,m\}:\;
  \fp\!\bigl(\sel[d \mapsto \Rep_d(C_i)]\bigr)
  = \fp\!\bigl(\sel[d \mapsto \Rep_d(C_j)]\bigr).
\]
\end{definition}

\begin{definition}[Superposition Set]\label{def:superset}
During the DFS, a set $\Sigma \subseteq \Dim$ tracks which dimensions
have been \emph{superimposed} (deferred). When $d$ is superimposable, the
DFS explores a single representative and adds $d$ to $\Sigma$,
rather than branching over all class representatives.
\end{definition}

\begin{definition}[Super Key]\label{def:superkey}
The \textbf{super key} generalizes the irrelevance key by additionally
wildcarding superimposed dimensions:
\[
  \mathsf{superKey}(\sel,\Sigma) = \bigl(k_1,\dots,k_n\bigr)
  \quad\text{where}\quad
  k_i =
  \begin{cases}
    \wildcard & \text{if } d_i \in \Sigma \text{ or }
                \Irr(d_i,\sel),\\
    c_i       & \text{otherwise}.
  \end{cases}
\]
Note that $\ik(\sel) = \mathsf{superKey}(\sel, \emptyset)$.
\end{definition}

\begin{definition}[Collapse]\label{def:collapse}
A superimposed dimension $d \in \Sigma$ must be \textbf{collapsed} at
state $\sel$ with next dimension $d^*$ if setting $\sel(d)$ to different
class representatives of $d$ would change either:
\begin{enumerate}[label=(\alph*)]
  \item $\mathsf{next}(\sel)$, or
  \item the set of enabled edges (by class) of $d^*$.
\end{enumerate}
Collapse replaces the single-path exploration with branching over all
class representatives of $d$, removing $d$ from $\Sigma$.
\end{definition}

\begin{algorithm}[H]
\caption{Superposition DFS}\label{alg:superdfs}
\begin{algorithmic}[1]
\Procedure{SuperDFS}{$\sel, \Sigma$}
  \State $\mathit{key} \gets \mathsf{superKey}(\sel, \Sigma)$
  \State $d^* \gets \mathsf{next}(\sel)$
  \State $E^* \gets \{e \in \Edg{d^*} : \Dis_{d^*,e}(\sel)=0\}$
    if $d^* \neq \bot$, else $\emptyset$
  \If{$d^* = \bot$ or $E^* = \emptyset$}
    \Comment{Terminal or dead end --- always compute fresh}
    \State \Call{RecordReach}{$\sel, \Sigma, 0$}
    \State \Return $M(\sel)$
  \EndIf
  \If{$\mathit{key} \in \mathit{visited}$}
    \Comment{Non-terminal memoization}
    \State \Return cached mask for $\mathit{key}$
  \EndIf
  \State $\mathit{visited} \gets \mathit{visited} \cup \{\mathit{key}\}$
  \ForAll{$d \in \Sigma$}
    \If{$\mathrm{NeedsCollapse}(d, \sel, d^*)$}
      \State $\mathit{mask} \gets 0$
      \ForAll{$r \in \Rep_d$} \Comment{branch on class reps of $d$}
        \State $\sel' \gets \clean(\sel[d \mapsto r])$
        \State $\mathit{mask} \gets \mathit{mask} \mathbin{|}
               \Call{SuperDFS}{\sel', \Sigma \setminus \{d\}}$
      \EndFor
      \State \Call{RecordReach}{$\sel, \Sigma, \mathit{mask}$}
      \State \Return $M(\sel) \mathbin{|} \mathit{mask}$
    \EndIf
  \EndFor
  \State Deduplicate $E^*$ to one edge per class
  \If{only one class is enabled}
    \State $\mathit{child} \gets \Call{SuperDFS}{\clean(\sel[d^* \mapsto e]),\;\Sigma}$
  \ElsIf{$d^*$ is superimposable at $\sel$}
    \State $\mathit{child} \gets \Call{SuperDFS}{\clean(\sel[d^* \mapsto \Rep_{d^*}(C_1)]),\;
           \Sigma \cup \{d^*\}}$
  \Else
    \State $\mathit{child} \gets 0$
    \ForAll{class representatives $r$ of enabled edges of $d^*$}
      \State $\mathit{child} \gets \mathit{child} \mathbin{|}
             \Call{SuperDFS}{\clean(\sel[d^* \mapsto r]),\;\Sigma}$
    \EndFor
  \EndIf
  \State $\mathit{mask} \gets M(\sel) \mathbin{|} \mathit{child}$
  \State cache $\mathit{mask}$ for $\mathit{key}$
  \State \Call{RecordReach}{$\sel, \Sigma, \mathit{child}$}
  \State \Return $\mathit{mask}$
\EndProcedure
\end{algorithmic}
\end{algorithm}

\begin{theorem}[Reduction R3: Superposition Correctness]\label{thm:super}
For all states $\sel$ and outcomes $k$:
\[
  \Reach_{\textsf{ClassDFS}}(\sel,k)
  = \Reach_{\textsf{SuperDFS}}(\sel,\emptyset,k).
\]
Superposition with deferred collapse produces the same reachability
answers, detects the same dead ends, and triggers the same validation
violations as the class-merged DFS without superposition.
\end{theorem}

\begin{remark}[Terminal Memoization Caveat]\label{rem:termmemo}
The super key $\mathsf{superKey}(\sel,\Sigma)$ wildcards both
superimposed and irrelevant dimensions. At terminal states (no remaining
visible unanswered dimensions), \emph{every} dimension satisfies the
irrelevance criterion---there are no future decisions to affect---so the
super key collapses to the all-wildcard string $(**\cdots*)$ regardless of
the actual state values.

However, different terminal states can match different outcome predicates
$O_k$. If the DFS memoizes by super key \emph{before} checking terminality,
the first terminal visited caches its outcome mask under the all-wildcard
key, and all subsequent terminals return that cached mask instead of
evaluating $M(\sel)$ against their own state. This produces both false
positives (outcomes attributed to unreachable terminals) and false
negatives (outcomes from later terminals never recorded).

The corrected algorithm (Algorithm~\ref{alg:superdfs}) checks terminality
\emph{before} consulting the memoization cache, ensuring each terminal
always computes its mask fresh. Non-terminal states still benefit from
super key memoization, since their reachability depends on subgraph
structure (captured by the super key) rather than on the specific state
values.
\end{remark}

%--------------------------------------------------------------------------
\subsection{Reachability Recording}\label{sec:recordreach}
%--------------------------------------------------------------------------

The DFS does not simply return masks up the call stack; it also writes
into a persistent map $\mathsf{ReachMap}$ keyed by irrelevance keys.
Multiple DFS paths may map to the same $\ik$, and their masks are
OR-accumulated.

\begin{definition}[RecordReach]\label{def:recordreach}
Given state $\sel$, superposition set $\Sigma$, and child mask
$\mathit{cm}$ (the OR of masks returned by recursive calls), the
\textsf{RecordReach} procedure performs three steps:
\begin{enumerate}[label=(\roman*)]
  \item \textbf{Direct recording.}
    Compute $\ik(\sel)$ and set
    \[
      \mathsf{ReachMap}(\ik(\sel)) \;\gets\;
      \mathsf{ReachMap}(\ik(\sel)) \;\mathbin{|}\;
      M(\sel) \;\mathbin{|}\; \mathit{cm}.
    \]
  \item \textbf{Superposition expansion.}
    If $\Sigma \neq \emptyset$, enumerate
    $\prod_{d \in \Sigma} \Edg{d}$ (all raw edge values for each
    superimposed dimension). For each assignment, temporarily override
    $\sel$, recompute $\ik$, and OR-accumulate
    $M(\sel') \mathbin{|} \mathit{cm}$ into the corresponding entry.
  \item \textbf{Variant expansion.}
    For each primary dimension $d \in \mathcal{P}$ where
    $\sel(d) \neq \bot$, enumerate all same-class, non-disabled edges.
    For each combination, recompute $\ik$ and OR-accumulate
    $M(\sel') \mathbin{|} \mathit{cm}$.
\end{enumerate}
\end{definition}

\begin{remark}[OR-Accumulation Semantics]
The $\mathsf{ReachMap}$ is monotonically increasing: each write can only
set additional bits, never clear them. This is correct because the map
stores the union of all outcomes reachable from any state that projects to
a given $\ik$. The OR-accumulation means the DFS need not track which
specific path contributed which bits---the final map entry is the union
over all contributing paths.
\end{remark}

%--------------------------------------------------------------------------
\subsection{Variant Expansion (\texttt{noClassMergeDims})}
\label{sec:variants}
%--------------------------------------------------------------------------

Equivalence class merging can lose information when outcome predicates
distinguish individual edge values within a class.

\begin{definition}[Primary Dimension]\label{def:primary}
A dimension $d$ is a \textbf{primary dimension} if some outcome variant
$O_{k,v}$ requires $\hat\sel(d) = v$, where $v$ and $v'$ with $v \neq v'$
may belong to the same equivalence class $[v]_d = [v']_d$.
\end{definition}

\begin{definition}[Variant Expansion]\label{def:expand}
For a set $\mathcal{P} \subseteq \Dim$ of primary dimensions, after the
DFS records a state $\sel$ with reachability mask $m$, the
\textbf{variant expansion} enumerates the Cartesian product:
\[
  \prod_{d \in \mathcal{P},\;\sel(d)\neq\bot}
  \{e \in \Edg{d} : [e]_d = [\sel(d)]_d
   \;\wedge\; \Dis_{d,e}(\sel)=0\}
\]
For each assignment in the product, $\ik$ is recomputed and the
reachability mask (including fresh $M(\cdot)$ checks) is OR-accumulated
into $\mathsf{ReachMap}$.
\end{definition}

\begin{theorem}[Reduction R4: Variant Expansion Correctness]\label{thm:variant}
Let $\textsf{ClassDFS}_\mathcal{P}$ denote the class-merged DFS that does
\emph{not} merge classes on dimensions in $\mathcal{P}$ (i.e., branches
on all enabled edges, not just representatives, for $d \in \mathcal{P}$).
Then for all states $\sel$ and outcomes $k$:
\[
  \Reach_{\textsf{ClassDFS}_\mathcal{P}}(\sel,k)
  = \Reach_{\textsf{SuperDFS}+\text{Expand}_\mathcal{P}}(\sel,k).
\]
Variant expansion at recording time produces the same per-edge
reachability answers as not merging those dimensions during traversal.
This is needed because outcome predicates can distinguish individual edge
values that the class-merged DFS would otherwise conflate.
\end{theorem}

%--------------------------------------------------------------------------
\subsection{Superposition Expansion}\label{sec:superexpand}
%--------------------------------------------------------------------------

Similarly, superimposed dimensions must be expanded when recording
reachability, since outcome predicates may depend on the deferred value.

\begin{definition}[Superposition Expansion]\label{def:superexpand}
When recording reachability for state $\sel$ with superposition set
$\Sigma$, enumerate all combinations of raw edge values for each
$d \in \Sigma$:
\[
  \prod_{d \in \Sigma} \Edg{d}
\]
For each combination, recompute $\ik$ and OR-accumulate
$M(\cdot)$ and child masks into $\mathsf{ReachMap}$.
\end{definition}

\begin{remark}
Superposition expansion iterates over \emph{all} raw edge values (not
just class representatives or enabled edges), because the DFS deferred
$d$'s choice entirely---any value could have been chosen at the original
branching point.
\end{remark}


%==========================================================================
\section{Supporting Lemmas}\label{sec:lemmas}
%==========================================================================

The reduction theorems in Section~\ref{sec:reductions} each depend on
supporting lemmas that establish the structural properties making the
reductions safe. Together with the theorems, these form the full set of
claims to verify: if these hold, the production traversal is guaranteed to
produce the same results as the baseline DFS for every algorithm.

\begin{lemma}[Clean Termination]\label{lem:term}
The clean operator $\clean(\sel)$ reaches a fixpoint in at most
$P \cdot |\Dim|$ individual dimension inspections, where $P = 5$.
\end{lemma}

\begin{proof}[Proof sketch]
Each pass iterates over $|\Dim|$ dimensions. A pass that makes no change
terminates the loop. Each change deletes a value; since there are
finitely many dimensions and deletions are monotone, the iteration
converges.
\end{proof}

\begin{lemma}[Class Congruence]\label{lem:congruence}
The equivalence class partition $\sim_d$ is a congruence with respect to
the activation, hiding, and disabling predicates: for all $d' \neq d$ and
all $e_1 \sim_d e_2$,
\[
  \Act_{d'}(\sel[d \mapsto e_1]) = \Act_{d'}(\sel[d \mapsto e_2]),\quad
  H_{d'}(\sel[d \mapsto e_1]) = H_{d'}(\sel[d \mapsto e_2]),
\]
\[
  \Dis_{d',e'}(\sel[d \mapsto e_1]) = \Dis_{d',e'}(\sel[d \mapsto e_2])
  \quad \forall\, e' \in \Edg{d'},
\]
\[
  \res{\sel[d \mapsto e_1]}{d'} = \res{\sel[d \mapsto e_2]}{d'}
  \quad \forall\, \text{derived } d'.
\]
\end{lemma}

\begin{lemma}[Irrelevance Well-Definedness]\label{lem:irr-wd}
$\Irr(d,\sel)$ depends only on $\ck(\sel)$ and whether $\sel(d) \neq
\bot$, not on the specific value of $\sel(d)$ within its class. Hence
$\ik(\sel)$ is well-defined: two states with the same class key produce
the same irrelevance vector.
\end{lemma}

\begin{lemma}[Superposition--Collapse Equivalence]\label{lem:super-collapse}
Let $d$ be superimposable at $\sel$ with enabled class set
$\{C_1,\dots,C_m\}$. Then for any outcome $k$:
\[
  \bigvee_{i=1}^{m}
  \Reach\bigl(\clean(\sel[d \mapsto \Rep_d(C_i)]),\, k\bigr)
  \;=\;
  \Reach\bigl(\clean(\sel[d \mapsto \Rep_d(C_1)]),\, k\bigr).
\]
If $d$ is superimposable, exploring a single representative suffices
because the fingerprint identity guarantees all branches encounter
identical subproblems.
\end{lemma}

\begin{lemma}[Terminal Memoization Soundness]\label{lem:termmemo}
Let $\sel_1, \sel_2$ be two terminal states with
$\mathsf{superKey}(\sel_1,\Sigma_1) = \mathsf{superKey}(\sel_2,\Sigma_2)$.
In general, $M(\sel_1) \neq M(\sel_2)$: the outcome masks may differ
because the outcome predicates $O_k$ evaluate against the actual resolved
state, not the super key. Therefore, memoizing $M(\sel)$ by super key at
terminal states is unsound.

At non-terminal states, the returned mask is
$M(\sel) \mathbin{|} \mathit{childMask}$, where
$\mathit{childMask}$ depends on the subgraph structure. Two non-terminal
states with the same super key have identical subgraph structure (by the
fingerprint and irrelevance invariants), so $\mathit{childMask}$ is the
same. Although $M(\sel_1)$ and $M(\sel_2)$ may also differ, the
\textsf{RecordReach} procedure writes the correct per-state mask into
$\mathsf{ReachMap}$ regardless of which cached $\mathit{childMask}$ is
used. Therefore, super key memoization at non-terminal states is sound
for reachability recording, provided terminals are always computed fresh.
\end{lemma}

\begin{remark}[Edge-Local Writes Replace Lazy Derivation]\label{rem:crossover}
Earlier revisions of the graph supported lazy
\texttt{deriveWhen} rules: a node $d_1$ could declare a derivation rule
that recomputed its effective value on demand from other dimensions
(e.g.\ ``if $d_2 = \mathsf{stall\_recovery}$ is $\mathsf{mild}$, then
$d_1$ resolves to $\mathsf{singularity}$''). Because the resolved value
depended on later answers, picking $d_2$ could retroactively change the
visibility of downstream dimensions, creating \textbf{path crossovers}
that complicated both the static analysis (which had to materialize
JIT Cartesian-product derivation tables) and the proof obligations on
class computation and fingerprint checks.

The current model eliminates this entirely. Every change to $\sel$ /
$\flav$ is committed by the edge that fires, either directly (via
\texttt{effects.set} / \texttt{setFlavor}) or via a module's
\texttt{exitPlan}. \texttt{resolvedVal} reduces to a direct
$\sel[d]$ lookup, and \texttt{matchCondition} reads $\sel[k]$ without
recursion. Dimensions that the user never picks (such as
$\mathsf{ruin\_type}$ or $\mathsf{governance}$) are written into
$\sel$ / $\flav$ by the upstream edges that own the triggering signal,
so the same crossover behavior is preserved as a plain forward push
rather than a lazy backward inference. This collapses the proof
obligations on class computation and fingerprint checks to those of an
ordinary monotone DAG.
\end{remark}

\begin{lemma}[Safe Push Optimization]\label{lem:safepush}
A dimension $d$ is a \textbf{safe push dimension} if $d$ does not appear
(directly or via derivation chains) in any edge-local condition
(\texttt{disabledWhen} or \texttt{requires}) of any other dimension.
For safe push dimensions, $\clean(\sel[d\mapsto e]) = \sel[d\mapsto e]$
--- clean selection has no effect, so it can be skipped.
\end{lemma}

\begin{proof}[Proof sketch]
If $d$ does not affect any edge disabling predicate, assigning $d$ cannot
invalidate any existing choice and cannot enable or disable edges
elsewhere. Thus clean selection finds nothing to change.
\end{proof}


%==========================================================================
\section{Empirical Validation}\label{sec:empirical}
%==========================================================================

To supplement formal verification, a test harness compares the optimized
traversal against a brute-force baseline DFS on a suite of small graphs.
The harness (\texttt{test/run.js}) works as follows:

\begin{enumerate}
  \item \textbf{Module injection.} Each test graph is a JSON file
    defining nodes, edges, conditions, and outcomes. The harness injects
    it into the engine and walker via Node.js module cache replacement,
    reusing the production code paths exactly.
  \item \textbf{Baseline DFS.} A brute-force recursive DFS explores every
    enabled edge at every visible non-derived dimension, with no class
    merging, irrelevance, or superposition. It records the outcome mask
    $M(\sel)$ at each visited state.
  \item \textbf{Optimized computation.} The production
    \textsf{computeReachability} function (Algorithm~\ref{alg:superdfs}
    with \textsf{RecordReach}) runs on the same graph and outcomes.
  \item \textbf{Per-state reachability comparison.} For every state
    $\sel$ visited by the baseline, the harness computes $\ik(\sel)$ and
    compares the true mask $M_\text{true}(\sel)$ against
    $\mathsf{ReachMap}(\ik(\sel))$:
    \begin{itemize}
      \item \textbf{No false positives} (\texttt{reach-fp}): every bit
        in $\mathsf{ReachMap}(\ik)$ must be set in
        $M_\text{true}(\sel)$.
      \item \textbf{No false negatives} (\texttt{reach-fn}): every bit
        in $M_\text{true}(\sel)$ must appear in
        $\mathsf{ReachMap}(\ik)$.
      \item \textbf{No missing keys} (\texttt{reach-missing}):
        every reachable $\ik$ whose true mask is non-zero must have an
        entry in $\mathsf{ReachMap}$ (the browser's
        \texttt{reachSet.has} query is fail-closed on absent keys).
    \end{itemize}
    This is strictly stronger than aggregating baseline masks by $\ik$
    before comparing: if two states with different true reach-sets
    collapse to the same $\ik$, the stored mask is their OR, which would
    appear correct to an aggregated check but violates the browser's
    per-state query.
  \item \textbf{Browser-sim comparison.} A separate harness
    (\texttt{test/reach-browser-sim.js}) replays the browser's
    \texttt{wouldReachOutcome} decision (lightweight push plus
    $\mathsf{ReachMap}$ lookup) at every reachable state and compares it
    against the true reach-set membership obtained by committing the push
    and descending in the baseline. Both \texttt{sim-fp} (UI green-lights
    a dead end) and \texttt{sim-fn} (UI hides a valid path) must be
    empty.
\end{enumerate}

\subsection{Test Graphs}

\begin{center}
\begin{tabular}{lrrrl}
\textbf{Graph} & \textbf{Dims} & \textbf{Baseline} & \textbf{irrKeys}
  & \textbf{Primary target} \\
\hline
\texttt{trivial}             & 3  & 15   & 4  & Sanity check \\
\texttt{classes}             & 3  & 23   & 4  & R1: equivalence classes \\
\texttt{irrelevance}         & 3  & 10   & 4  & R2: irrelevance wildcarding \\
\texttt{superposition}       & 4  & 36   & 5  & R3: superposition \\
\texttt{deep-superposition}  & 5  & 118  & 10 & R3: chained superposition \\
\texttt{variants}            & 3  & 4    & 4  & R4: variant expansion \\
\texttt{kitchen-sink}        & 7  & 23   & 13 & Combined features \\
\texttt{edge-write-crossover} & 5  & 9    & 7  & Edge-local write crossover \\
\texttt{chained-edge-writes}  & 5  & 7    & 6  & Transitive edge writes \\
\texttt{convergence-basic}   & 7  & 14   & 9  & Path convergence \\
\texttt{diamond-convergence} & 4  & 8    & 7  & Diamond-shaped subgraph \\
\texttt{disabled-by-derived} & 4  & 8    & 7  & Derived-value disabling \\
\texttt{edge-requires}       & 4  & 16   & 12 & \texttt{requires} clauses \\
\texttt{fp-cascade}          & 7  & 3    & 3  & Fingerprint cascade \\
\texttt{hide-when}           & 4  & 12   & 12 & \texttt{hideWhen} predicates \\
\texttt{multi-activate}      & 6  & 35   & 15 & Disjunctive activation \\
\texttt{outcome-only-dim}    & 4  & 5    & 5  & Outcome-only dimensions \\
\texttt{rederivation}        & 6  & 36   & 27 & Re-derived dead-end check \\
\texttt{production-subgraph} & 16 & 1204 & 20
  & Real graph (capability $\to$ open\_source) \\
\end{tabular}
\end{center}

The \texttt{edge-write-crossover} graph is the minimal reproducer for
the terminal memoization bug (Remark~\ref{rem:termmemo}): 5 dimensions,
10 edges, with an upstream edge whose \texttt{effects.set}
write changes a downstream dimension's effective value, creating a
path crossover. Before the fix (moving the terminal check before
memoization), it produced both false positives and false negatives.
After the fix, all 19 graphs pass.

The \texttt{production-subgraph} graph contains the first 16 dimensions
of the production Singularity Map (up to and including
\texttt{open\_source}), with all real \texttt{activateWhen},
\texttt{disabledWhen}, \texttt{requires}, and
\texttt{effects.when} conditions preserved. Its 1204 baseline states compress to 20 irrelevance
keys, and the optimized--baseline comparison passes with zero
discrepancies (no reachability false positives, false negatives, or
missing entries; browser-sim FP/FN both zero).

\subsection{Limitations}

On small graphs, the iterative class computation often collapses to a
single equivalence class per dimension (insufficient structural depth to
sustain multi-class behavior). This makes the R1 and R3 tests pass
trivially. The \texttt{production-subgraph} provides meaningful
multi-class testing: \texttt{capability} (2 classes),
\texttt{stall\_recovery} (2 classes).

The baseline DFS does not scale to the full production graph
(${\sim}216\text{k}$ states visited by the walk, ${\sim}243\text{k}$
irrelevance keys in the final $\mathsf{ReachMap}$) within reasonable
time, so the full graph is validated only by the \textsf{walk}
procedure's internal invariant checks and the browser-sim harness, not
by baseline comparison.


%==========================================================================
\section{Summary of Claims for Verification}\label{sec:summary}
%==========================================================================

The central verification goal is:

\begin{quote}
\emph{The production traversal (SuperDFS + variant/superposition
expansion + $\ik$-memoization) produces exactly the same results as the
baseline DFS for every algorithm defined in Section~\ref{sec:algorithms}
and every validation invariant defined in Section~\ref{sec:validate}.}
\end{quote}

\noindent
This decomposes into the four reduction theorems (R1--R4) and seven
supporting lemmas listed below. The validation invariants define what
``same results'' means concretely. Section~\ref{sec:empirical} provides
empirical evidence via baseline-vs-optimized comparison on 8~test graphs.

\bigskip

\begin{center}
\begin{tabular}{llp{7.5cm}}
\textbf{Ref} & \textbf{Type} & \textbf{Statement} \\
\hline
\multicolumn{3}{l}{\emph{Reduction theorems
  (does the reduced DFS match the baseline?)}} \\[2pt]
\ref{thm:class}
  & Thm (R1)
  & Class-merged DFS preserves all algorithm results. \\
\ref{thm:irr}
  & Thm (R2)
  & $\ik$-memoization preserves all algorithm results. \\
\ref{thm:super}
  & Thm (R3)
  & Superposition DFS preserves all algorithm results. \\
\ref{thm:variant}
  & Thm (R4)
  & Variant expansion restores per-edge reachability. \\[4pt]
\hline
\\[-6pt]
\multicolumn{3}{l}{\emph{Supporting lemmas
  (structural properties that make the reductions safe)}} \\[2pt]
\ref{lem:clean-term}
  & Lemma
  & Clean selection terminates in $O(P \cdot |\Dim|)$. \\
\ref{lem:congruence}
  & Lemma
  & Equivalence classes form a congruence over predicates. \\
\ref{lem:irr-wd}
  & Lemma
  & Irrelevance key is well-defined (depends only on class key). \\
\ref{lem:super-collapse}
  & Lemma
  & Superposition of fingerprint-identical branches is safe. \\
\ref{lem:termmemo}
  & Lemma
  & Super key memoization is sound only at non-terminals. \\
\ref{lem:safepush}
  & Lemma
  & Safe push dimensions can skip clean selection. \\[4pt]
\hline
\\[-6pt]
\multicolumn{3}{l}{\emph{Algorithms and validation invariants
  (what ``same results'' means)}} \\[2pt]
\ref{def:reach}
  & Defn
  & Reachability: $\exists$ play from $\sel$ to terminal matching $O_k$. \\
\ref{def:deadend}
  & Defn
  & Dead ends: no outcome, no enabled next edge. \\
\ref{def:edgecov}
  & Defn
  & Edge coverage: fraction of $(d,e)$ pairs visited. \\
\ref{def:deaded}
  & Defn
  & Statically dead edges are never enabled. \\
\ref{def:rederived}
  & Defn
  & Re-derived dead ends: derivation hides a question. \\
\ref{def:stuck}
  & Defn
  & Stuck nodes: visible dimension with all edges disabled. \\
\ref{def:ambig}
  & Defn
  & Ambiguous outcomes: multiple templates match one state. \\
\ref{def:stackint}
  & Defn
  & Stack integrity: display order matches stack order. \\
\end{tabular}
\end{center}


\end{document}
